\documentclass[12pt]{article}
\usepackage[T1]{fontenc}
\usepackage[utf8]{inputenc}
\usepackage{lmodern}
\usepackage{textcomp}
\usepackage{enumitem}
\usepackage{geometry}
\geometry{margin=1in}
\begin{document}
\begin{center}
Answer Key - Philosophy - Undergraduate Level Assessment 25 \
\small (Answers are ordered exactly as in the question paper.)
\end{center}

\section*{Multiple Choice}
\begin{enumerate}[label=\arabic*.]
\item \textbf{Answer:} C
\\ \textit{Options:} A) 13; B) 18; C) 24; D) 60; E) 5
\\ \textit{Note:} The correct answer is 24 because in Jain philosophy, it is believed that 24 Tirthankaras, or Jinas, have attained enlightenment and provided spiritual guidance during the current time cycle known as the present age, or "Kali Yuga."
\item \textbf{Answer:} A
\\ \textit{Options:} A) preference.; B) natural law.; C) self-interest.; D) maximizing utility.; E) ethical relativism.
\\ \textit{Note:} Anscombe argues that moral obligations arise from individual preferences, as they shape the decisions and actions that reflect one's ethical commitments, emphasizing that moral judgments are closely tied to personal values and desires.
\item \textbf{Answer:} D
\\ \textit{Options:} A) this policy is morally problematic because it disrupts the natural order.; B) this policy is morally problematic because it negatively impacts biodiversity.; C) this policy is morally acceptable because it does not harm human beings.; D) this policy is morally unproblematic.; E) this policy is morally obligatory because it prioritizes human beings.
\\ \textit{Note:} Baxter's ethical framework prioritizes human interests and welfare, suggesting that if a policy does not harm humans, the extinction of animal species can be deemed morally acceptable or unproblematic.
\item \textbf{Answer:} C
\\ \textit{Options:} A) a capacity for emotional intelligence; B) a capacity for scientific understanding; C) a capacity for communal or friendly relationships; D) a capacity for certain kinds of feelings or sentiments; E) a capacity for independent thinking
\\ \textit{Note:} Metz argues that dignity is rooted in the ability to engage in communal or friendly relationships because these interactions highlight the intrinsic value of individuals and their capacity for mutual respect and connection, which is essential for recognizing and upholding human dignity.
\item \textbf{Answer:} A
\\ \textit{Options:} A) how the decision to preserve the environment benefits the environment.; B) how the destruction of the environment affects the economy.; C) why people who preserve the environment might be good people.; D) how destroying the environment affects future generations.; E) why the act of destroying nature might be immoral.
\\ \textit{Note:} Hill shifts the focus from the morality of destruction to the positive impacts of preservation, emphasizing the importance of understanding the benefits of environmental conservation rather than solely critiquing harmful actions.
\end{enumerate}

\section*{True / False}
\begin{enumerate}[label=\arabic*.]
\setcounter{enumi}{5}
\item \textbf{Answer:} False
\\ \textit{Options:} A) True; B) False
\\ \textit{Note:} Norcross argues that moral agency is contingent upon certain cognitive abilities and reasoning capabilities, which means not all beings are considered moral agents.
\item \textbf{Answer:} False
\\ \textit{Options:} A) True; B) False
\\ \textit{Note:} The argument that the universe must have a maker does not necessarily commit the fallacy of composition, as it draws on the principle that complex entities can have causes independent of the properties of their individual components.
\item \textbf{Answer:} False
\\ \textit{Options:} A) True; B) False
\\ \textit{Note:} Velleman argues that offering euthanasia options may undermine patient autonomy by framing the choice in a way that pressures individuals into opting for death rather than enhancing their treatment choices.
\item \textbf{Answer:} True
\\ \textit{Options:} A) True; B) False
\\ \textit{Note:} The statement is true because the predicate logic expression 'Bhgh' correctly represents the relationship of borrowing between George and Hector's lawnmower, with 'B' denoting the borrowing action and the appropriate arguments for George and the lawnmower.
\item \textbf{Answer:} True
\\ \textit{Options:} A) True; B) False
\\ \textit{Note:} Euthanasia is accurately defined as the intentional act of terminating a person's life to alleviate unbearable suffering, reflecting its primary ethical goal of compassion and relief from pain.
\end{enumerate}

\section*{Long Form Response}
\begin{enumerate}[label=\arabic*.]
\setcounter{enumi}{10}
\item \textbf{Answer:} The imposition of the death penalty raises significant moral concerns, particularly regarding the irreversible nature of executing innocent individuals. The possibility of wrongful convictions is an inherent flaw in the justice system, where errors can lead to the loss of innocent lives, thereby undermining the sanctity of human life. This reality not only inflicts profound harm on the victims' families but also casts a shadow over the moral authority of the state to enact such a final punishment. Furthermore, the death penalty may foster a societal desensitization to violence, raising ethical questions about the value placed on life and the potential for a culture that endorses retribution over rehabilitation. Ultimately, the risk of executing the innocent challenges the legitimacy of capital punishment and compels us to reconsider its moral implications within a just society.
\\ \textit{Note:} This answer is correct because it effectively highlights the moral implications of the death penalty, particularly the irreversible harm caused by executing innocent individuals, which challenges the ethical legitimacy of state-sanctioned capital punishment.
\item \textbf{Answer:} In Jaina traditions, the concept of ajiva refers to "non-soul," encompassing all that is non-living and devoid of consciousness. Ajiva contrasts with jiva, or the soul, which is seen as eternal and sentient. This distinction is crucial in Jaina philosophy and metaphysics, as it underscores the belief that only jiva is capable of experiencing liberation through spiritual advancement, while ajiva represents the material world that binds souls in cycles of birth and death. The understanding of ajiva fosters a deep respect for all forms of life and non-violence, as Jains strive to minimize harm to any living being, recognizing the interconnectedness of existence. Thus, ajiva not only shapes the metaphysical landscape of Jainism but also informs ethical practices and the pursuit of spiritual purity.
\\ \textit{Note:} correct because it clearly distinguishes between ajiva and jiva, emphasizing that ajiva represents the non-conscious, material aspects of existence, while jiva signifies the conscious soul capable of spiritual liberation, which is central to Jaina metaphysics.
\end{enumerate}

\vfill
\noindent\textit{End of Answer Key}
\end{document}